%%%%%%%%%%%%%%%%%%%%%%%%%%%%%%%%%%%%%%%%%%%%%%%%%%%%%%%%%%%%%%%%%%%%%%%%%%%%%%%%
%   Copyright 2009, Oracle and/or its affiliates.
%   All rights reserved.
%
%
%   Use is subject to license terms.
%
%   This distribution may include materials developed by third parties.
%
%%%%%%%%%%%%%%%%%%%%%%%%%%%%%%%%%%%%%%%%%%%%%%%%%%%%%%%%%%%%%%%%%%%%%%%%%%%%%%%%

\section{Static Expressions}
\seclabel{static-expr}

\note{Of the static expressions in a size, the interpreter computes a sum, a
difference, a product and a power exactly, and refuses a power at a negative
exponent; the compiled type checker computes each operation on numerals,
refusing as out of range a power it does not compute, and compares other size
expressions by their written form (\chapref{types}); the compiled path
computes a size when it names a class. The compiled type checker does not
infer a \KWD{bool}, \KWD{dim} or \KWD{unit} static argument.}

\emph{Static expressions} denote \emph{static values}.  Conceptually,
a static expression is not evaluated while a program is running, and
thus its evaluation
cannot complete abruptly.  Given instantiations of all static
parameters (described in \chapref{trait-parameters}) in scope of a
static expression, the value of the static expression can be
determined statically.  Static expressions which use a restricted subset of operations can occur as
\emph{static arguments} which instantiate static parameters and
\KWD{where}-clause constraints (described in \secref{where-clauses}).
We define the set of static
expressions by first defining the types of static expressions, and
distinguishing static values from the closely related literal values.
We then describe the expressions that evaluate to the various kinds of
static values.


\subsection{Types of Static Expressions}

There are three groups of traits that describe literals
and static expressions:
\begin{enumerate}
\item The \emph{literal} traits, which describe boolean,
  character, string, and numerals.  For example, the literal \VAR{true} has
  trait \EXP{\TYP{BooleanLiteral}\llbracket\VAR{true}\rrbracket} and a
  character literal has trait \TYP{Character}.
See \secref{literals} for a discussion of Fortress literals.

\item The \emph{constant} traits, which describe values denoted by
expressions composed from literals and operators (described in
\secref{constant-values});
the type of a constant expression encodes the value of the expression.
For example, the type of a constant expression \EXP{3+5},
%IntegerConstant[\false, 3+5\]
\EXP{\TYP{IntegerConstant}\llbracket\VAR{false}, 3+5\rrbracket},
encodes the value of the expression; the operator \EXP{+} is
used in the static arguments.  When an operator is not supported in static
arguments, \library\ can still encode the value of the expression by
providing a specific constant trait such as
\EXP{\TYP{RationalValueTimesPi}\llbracket\VAR{false},1,1\rrbracket}
(described in \secref{pi}).



\item The \emph{static} traits, which describe values denoted by
expressions composed from literals, operators (described in
\secref{constant-values}), and \KWD{bool}, \KWD{nat},
and \KWD{int} parameters; here the type
does not encode the value of the expression, but the value of the
expression can nevertheless be known statically if specific
values are specified for the static parameters.  Also, in situations
where the type of an expression composed solely from literals and
operators nevertheless cannot be described by a constant trait, then a
static trait may be used to describe it instead.
For example, a static expression \EXP{2(3+m)} where \VAR{m} is a
\KWD{nat} parameter has trait \TYP{NaturalStatic}.
\end{enumerate}

\Library\ declare none of the constant traits.
An expression composed only of numerals and operators is folded, before the
static check and on both implementations, into a \emph{signed constant}, which
holds its value exactly and whose kind its operands give: computed from
numerals without a radix point and without a quotient, it is an integer,
written as an integer numeral or a negated one; computed from numerals without
a radix point with a quotient among its operations, it is a quotient, written as
the reduced quotient of two integer numerals; and computed from a numeral with a
radix point, it is written as a numeral with a radix point, or a negated one,
where its denominator divides a power of ten, and otherwise as the reduced
quotient of two numerals whose numerator has a radix point.
Where nothing fixes the type of a signed constant (\secref{literals}), an
integer takes the default of an integer numeral, a quotient the type
\EXP{\mathbb{Q}}, and a constant computed from a numeral with a radix point the
type \EXP{\mathbb{R}64}, whatever its value, so that \EXP{6/3} is a
\EXP{\mathbb{Q}} and \EXP{0.5 + 0.5} an \EXP{\mathbb{R}64}; where a place fixes
a type, the constant is converted into it as a numeral is, with at most one
rounding.
So \EXP{2147483647 + 1} is the \EXP{\mathbb{Z}64} \EXP{2147483648},
\EXP{-2147483648} is a \EXP{\mathbb{Z}32}, \EXP{x = 7/2} binds the
\EXP{\mathbb{Q}} \EXP{7/2}, and \EXP{0.1 + 0.2 = 0.3} is \VAR{true}.
The operators whose result depends on the width of a type, the wrapping and
saturating operators, the bitwise operators and the shifts, are not folded, nor
is a division by zero, nor an operation whose floating-point result could be a
negative zero: their operands take their types, and the operator its meaning
there.
\revision{revival-numerals}{The Working Draft of February 2011 gave a constant
expression a constant trait whose static arguments encode its value, as
\EXP{\TYP{IntegerConstant}\llbracket\VAR{false}, 3+5\rrbracket}, and no library
declares those traits.
The revival decided on 8~November 2026 that an expression of numerals alone is
folded exactly into a signed constant, which takes one of three defaults where
nothing fixes its type: the value that a constant trait would encode in its
static arguments is held in the constant itself, written as numerals, so that a
coercion of a program's own sees only the four numbers of a numeral as written
(\secref{literals}).}


The only operation on literals that produces a new literal (as opposed
to a constant) is concatenation by using the operator \EXP{\|}.
One may concatenate mixed string and
character literals, producing a string literal.
For example, ``\txt{foo}''\EXP{\|}''\txt{bar}'' yields ``\txt{foobar}''.
One may also concatenate two natural numerals
of the same radix, producing a new natural numeral of that radix.
For example, \EXP{\mathtt{deadc}_{16} \| \mathtt{0de}_{16}} yields \EXP{\mathtt{deadc0de}_{16}}.
\revision{revival-string-literals}{The Working Draft of February 2011 states
the concatenation of string literals, which neither of the team's
implementations built.
The revival decided on 8~November 2026 that a concatenation of two string
literals, as of a string and a character literal, is folded into one string
literal before the static check, on both implementations, and so denotes one
object (\secref{literals}).}
\revision{revival-numerals}{The concatenation of two natural numerals of the
same radix is built as written: before the static check, on both
implementations, it is folded into one natural numeral, whose digits are the
digits of both numerals in order, leading zeros among them, so that its four
numbers (\secref{literals}) count them.}

Every literal trait extends an appropriate constant trait, and every
constant trait extends an appropriate static trait.  So every
literal is also a constant expression, and every constant expression
is a static expression.




\subsection{Static Expressions and Values}
\seclabel{constant-values}

Static parameters are static expressions.
A \KWD{nat} parameter denotes a value that has type \TYP{NaturalStatic}
(which extends \TYP{IntegerStatic}).
An \KWD{int} parameter denotes a value that has type \TYP{IntegerStatic}.
A \KWD{bool} parameter denotes a value that has type \TYP{BooleanStatic}.


Boolean static expressions may be combined using the operators
\EXP{\wedge}, \EXP{\vee}, \EXP{\oplus},
\EXP{\equiv}, \EXP{\leftrightarrow}, \OPR{NAND}, \OPR{NOR}, \EXP{=},
\EXP{\neq}, and \EXP{\rightarrow} (See \appref{operator-precedence} for a
discussion of Fortress operators.)
to produce other static expressions denoting boolean static expressions.
If both operands are boolean constant expressions, then the result is also
a boolean constant expression.

Character static expressions may be compared using the operators \EXP{<},
\EXP{\leq}, \EXP{\geq}, \EXP{>},
\EXP{=}, and \EXP{\neq} to produce boolean static expressions.  If both
operands are character constant expressions, then the result is a boolean
constant expression.

Numeric static expressions may be compared using the operators \EXP{<},
\EXP{\leq}, \EXP{\geq}, \EXP{>},
\EXP{=}, and \EXP{\neq} to produce boolean static expressions.  If both
operands are numeric constant expressions, then the result is a boolean
constant expression.

Numeric static expressions may be combined using the operators and
functions \EXP{+},
\EXP{-}, \EXP{\times}, \EXP{\cdot}, \EXP{/}, \EXP{!}, \OPR{MIN}, \OPR{MAX},
\EXP{\surd}, \VAR{floor}, \VAR{ceiling},
\VAR{hyperfloor}, \VAR{hyperceiling}, \EXP{\gcd}, \TYP{lcm}, \EXP{\sin},
\EXP{\cos}, \EXP{\tan}, \EXP{\arcsin}, \EXP{\arccos}, and \EXP{\arctan}
to produce new numeric static expressions.  If the result is indeed
a numeric static expression, and both operands are numeric constant
expressions, then the result is also a numeric constant expression
as long as the constant traits provided by \library\ can encode the value
of the expression.
